\section{Discussion}
\label{sec:discussion}

\paragraph{Why the loose prompt and not the vague one?}
The two prompts that produced no drift have something in common that the
loose prompt lacks: each one asks for a fix. The scoped prompt names it, and
the vague prompt describes a symptom and asks for it to go away. The loose
prompt asks for the code to be put ``into good shape'', which is a request
for improvement with no stated limit, and the rabbit holes are exactly the
things an agent can point to as improvements. Our reading is that agents
treat the request's \emph{register} as a statement of scope: in this ladder,
``get things into good shape'' led many agents to edit the planted helpers,
while requests for a fix, specific or not, did not. This is consistent with work on
sycophancy \citep{sharma2024towards,perez2023discovering}, where models adapt
their output to what they take the user to want, with the finding of
\citet{qu2026overeager} that removing an explicit scope statement from the
prompt raises the rate of out-of-scope actions by coding agents, and with the
widened diffs \citet{weinberger2026promptinduced} report for a cleanup
clause. It runs
against a common intuition, which we held when designing the study, that
underspecified prompts are the dangerous ones, and against the result of
\citet{ji2026coding} for agents working in a simulated operations
environment. In our data, vague runs used more tool calls and made no planted-function
edits, although they did make other out-of-scope edits, mostly to tests, in
about an eighth of runs.

We should be careful about how far this goes. Our loose prompts are one
template, and a different phrasing (``fix this and tidy up anything that
looks wrong'') would almost certainly change the rate. The claim we can
support is narrower: holding everything else fixed, a cleanup-flavored
request is enough to move most agents from no drift to frequent drift, and
an equally short request for a fix is not.

\paragraph{What pass/fail scoring cannot see.}
Most visible/hidden designs, including ours, were built on the assumption
that intent violations show up as behavior changes. In our data they mostly
do not. Agents that rewrote helpers usually preserved the behaviors the
hidden tests check, about a quarter of planted-function edits changed only
comments or docstrings, and the only trace of the excursion was in the diff.
This matters for how coding agents are evaluated. A benchmark that reports
only resolved rates, even with strong held-out tests, will rate an agent that
rewrites four functions it was not asked about the same as one that changes
one line. Developers do not experience those two agents as equivalent. We
think a scope measure belongs next to every pass rate, and that a
function-level diff against a list of out-of-scope symbols is a cheap way to
get one.

\paragraph{Practical advice.}
In these tasks, symptom-only requests did not lead agents to change the
planted helpers, and cleanup requests often did. Phrases like ``clean up'', ``get this into
shape'', or ``take a look at the related helpers'' are read by most current
agents as permission to change more than the bug. Among the arms we tested,
model choice mattered as well: the two Anthropic models held scope under the
same loose prompt that most other arms acted on.
